% ============================================================
% presentacion_coordinadores.tex
% Sistema de Análisis Institucional DCBI 2026
% Presentación ejecutiva para coordinadores de licenciatura
% UAM-Azcapotzalco · División de Ciencias Básicas e Ingeniería
% Dr. César Simón López-Monsalvo · junio 2026
% ============================================================
\documentclass[aspectratio=169,12pt]{beamer}

% ── Codificación y lengua ────────────────────────────────────
\usepackage[utf8]{inputenc}
\usepackage[T1]{fontenc}
\usepackage[spanish,mexico]{babel}

% ── Tipografía: Helvetica / sans-serif ───────────────────────
\usepackage{helvet}
\renewcommand{\familydefault}{\sfdefault}

% ── Paquetes de apoyo ────────────────────────────────────────
\usepackage{graphicx}
\usepackage{tikz}
\usetikzlibrary{arrows.meta, positioning, calc}
\usepackage{xcolor}
\usepackage{booktabs}
\usepackage{microtype}
\usepackage{array}

% ── Tema base: metropolis ────────────────────────────────────
\usetheme{metropolis}

% ── Paleta nowhere_minimal / dashboard DCBI ─────────────────
\definecolor{ink}{HTML}{1A1A1A}
\definecolor{bg}{HTML}{F5F4F1}
\definecolor{uamred}{HTML}{C82D23}
\definecolor{graymid}{HTML}{6B6B6B}
\definecolor{cardwhite}{HTML}{FFFFFF}
\definecolor{cardborder}{HTML}{E4E2DF}

% ── Configuración de colores Beamer ──────────────────────────
\setbeamercolor{normal text}{fg=ink, bg=bg}
\setbeamercolor{frametitle}{fg=bg, bg=ink}
\setbeamercolor{title separator}{fg=uamred}
\setbeamercolor{progress bar}{fg=uamred, bg=cardborder}
\setbeamercolor{structure}{fg=uamred}
\setbeamercolor{alerted text}{fg=uamred}
\setbeamercolor{block title}{fg=cardwhite, bg=uamred}
\setbeamercolor{block body}{fg=ink, bg=uamred!8}
\setbeamercolor{section in toc}{fg=ink}
\setbeamercolor{subsection in toc}{fg=graymid}

% ── Supresión de barras de navegación ────────────────────────
\setbeamertemplate{navigation symbols}{}
\metroset{numbering=none, progressbar=none, block=fill}

% ── Fuentes Beamer ────────────────────────────────────────────
\setbeamerfont{title}{size=\LARGE, series=\bfseries}
\setbeamerfont{subtitle}{size=\large, series=\normalfont}
\setbeamerfont{frametitle}{size=\large, series=\bfseries}
\setbeamerfont{author}{size=\small}
\setbeamerfont{date}{size=\small}
\setbeamerfont{institute}{size=\footnotesize}

% ── Comando KPI ───────────────────────────────────────────────
% Uso: \kpi{número}{etiqueta}
% Nota: no usar \\ en los argumentos — se separan en dos líneas
% por el salto de línea del nodo con align=center.
\newcommand{\kpi}[2]{%
  \begin{tikzpicture}[baseline]
    \node[
      fill=uamred,
      text=cardwhite,
      inner xsep=8pt,
      inner ysep=5pt,
      rounded corners=2pt,
      align=center,
      font=\sffamily\footnotesize\bfseries
    ] {#1\\[1pt]{\tiny\mdseries #2}};
  \end{tikzpicture}%
}

% =============================================================
\begin{document}
% =============================================================

%──────────────────────────────────────────────────────────────
% SLIDE 1 — PORTADA (fondo oscuro)
%──────────────────────────────────────────────────────────────
{
\setbeamercolor{background canvas}{bg=ink}
\begin{frame}[plain]
  \begin{tikzpicture}[remember picture, overlay]
    \fill[ink] (current page.south west) rectangle (current page.north east);

%    % Logo UAM-A — esquina superior izquierda
%    \node[anchor=north west, inner sep=0pt]
%      at ($(current page.north west)+(0.6cm,-0.5cm)$) {
%        \includegraphics[
%          width=4.2cm,
%          trim={48bp 64bp 473bp 35bp},
%          clip
%        ]{/Users/nowhere/Claude_Projects/Personal/logo_UAM.png}
%    };

    % Línea de acento roja — separador horizontal
    \draw[uamred, line width=1.2pt]
      ($(current page.north west)+(0.6cm,-1.2cm)$) --
      ($(current page.north east)+(-0.6cm,-1.2cm)$);

    % Título principal
    \node[anchor=north west, text=cardwhite,
          font=\sffamily\bfseries\LARGE, text width=13.5cm,
          inner sep=0pt]
      at ($(current page.north west)+(0.7cm,-1.7cm)$) {
        \hyphenpenalty=10000\exhyphenpenalty=10000
        Sistema de Análisis Institucional
    };

    % Subtítulo
    \node[anchor=north west, text=cardborder,
          font=\sffamily\large, text width=13.5cm,
          inner sep=0pt]
      at ($(current page.north west)+(0.7cm,-2.6cm)$) {
        División de Ciencias Básicas e Ingeniería · 2026
    };

    % Descriptor menor
    \node[anchor=north west, text=graymid,
          font=\sffamily\footnotesize, text width=13.5cm,
          inner sep=0pt]
      at ($(current page.north west)+(0.7cm,-3.4cm)$) {
        Cuatro proyectos de análisis curricular con inteligencia artificial
    };

    % Línea inferior
    \draw[uamred, line width=0.6pt]
      ($(current page.south west)+(0.6cm,2.5cm)$) --
      ($(current page.south east)+(-0.6cm,2.5cm)$);

    % Autor y fecha
    \node[anchor=south west, text=cardborder,
          font=\sffamily\footnotesize, text width=12cm, inner sep=0pt]
      at ($(current page.south west)+(0.7cm,1.5cm)$) {
        Dr.\ César Simón López-Monsalvo 
    };

    % QR — esquina inferior derecha
    \node[anchor=south east, inner sep=0pt]
      at ($(current page.south east)+(-0.7cm,0.6cm)$) {
        \includegraphics[width=2.5cm]{%
          /Users/nowhere/Claude_Projects/Admin_duties/dashboard/figuras/qr_dashboard_dcbi.png}
    };

    % URL bajo QR
    \node[anchor=north east, text=graymid,
          font=\sffamily\tiny, inner sep=0pt]
      at ($(current page.south east)+(-0.7cm,0.5cm)$) {
        cslmuam.github.io/dashboard-dcbi-2026
    };

  \end{tikzpicture}
\end{frame}
}

%──────────────────────────────────────────────────────────────
% SLIDE 2 — PANORAMA GENERAL (fondo claro)
%──────────────────────────────────────────────────────────────
{
\setbeamercolor{background canvas}{bg=bg}
\begin{frame}{Cuatro proyectos, una visión integrada}
  \vspace{0.15cm}

  % Diagrama de flujo — escala reducida para caber en 16:9
  \begin{center}
  \begin{tikzpicture}[
    scale=0.88, transform shape,
    box/.style={
      rectangle, fill=uamred, text=cardwhite,
      rounded corners=3pt, inner sep=7pt,
      font=\sffamily\small\bfseries,
      text width=2.9cm, align=center,
      minimum height=1.3cm
    },
    arr/.style={-{Stealth[length=5pt]}, line width=1.2pt, color=graymid}
  ]
    \node[box] (p1) at (0,0)   {CAT 26P\\[2pt]{\tiny Asignación de grupos}};
    \node[box] (p2) at (3.6,0) {Eficiencia\\terminal\\[2pt]{\tiny Diagnóstico de rezago}};
    \node[box] (p3) at (7.2,0) {Distribución\\docente\\[2pt]{\tiny Concentración y riesgo}};
    \node[box] (p4) at (10.8,0){Modificaciones\\2026\\[2pt]{\tiny Respuesta curricular}};

    \draw[arr] (p1.east) -- (p2.west);
    \draw[arr] (p2.east) -- (p3.west);
    \draw[arr] (p3.east) -- (p4.west);
  \end{tikzpicture}
  \end{center}

%  \vspace{0.35cm}
%  \begin{center}
%  \begin{tikzpicture}
%    \node[
%      fill=cardwhite, draw=cardborder, line width=0.6pt,
%      rounded corners=2pt, inner sep=9pt,
%      text width=13cm, align=center,
%      font=\sffamily\small\color{graymid}
%    ] {
%      Los cuatro proyectos se retroalimentan: el diagnóstico de los proyectos~2
%      y~3 informa las decisiones del proyecto~1 y las propuestas del proyecto~4.
%    };
%  \end{tikzpicture}
%  \end{center}
\end{frame}
}

%──────────────────────────────────────────────────────────────
% SLIDE 3 — DIVISOR SECCIÓN 01 (fondo oscuro)
%──────────────────────────────────────────────────────────────
{
\setbeamercolor{background canvas}{bg=ink}
\begin{frame}[plain]
  \begin{tikzpicture}[remember picture, overlay]
    \fill[ink] (current page.south west) rectangle (current page.north east);

    % Número grande semitransparente
    \node[anchor=west, text=uamred, opacity=0.18,
          font=\sffamily\bfseries\fontsize{160}{160}\selectfont,
          inner sep=0pt]
      at ($(current page.west)+(-1.0cm,0.0cm)$) {01};

    % Línea vertical de acento
    \draw[uamred, line width=2pt]
      ($(current page.center)+(-3.0cm,1.2cm)$) --
      ($(current page.center)+(-3.0cm,-1.2cm)$);

    % Título de sección
    \node[anchor=west, text=cardwhite,
          font=\sffamily\bfseries\LARGE, inner sep=0pt]
      at ($(current page.center)+(-2.6cm,0.38cm)$) {CAT 26P};

    % Subtítulo
    \node[anchor=west, text=cardborder,
          font=\sffamily\normalsize, inner sep=0pt]
      at ($(current page.center)+(-2.6cm,-0.38cm)$) {Asignación de grupos con inteligencia artificial};
  \end{tikzpicture}
\end{frame}
}

%──────────────────────────────────────────────────────────────
% SLIDE 4 — CAT 26P: PROBLEMA Y SOLUCIÓN (fondo claro)
%──────────────────────────────────────────────────────────────
{
\setbeamercolor{background canvas}{bg=bg}
\begin{frame}{CAT 26P: el problema y la solución}
  \begin{columns}[T]
    \begin{column}{0.48\textwidth}
      \begin{tikzpicture}
        \node[
          fill=cardwhite, draw=cardborder, line width=0.6pt,
          rounded corners=2pt, inner sep=9pt,
          text width=6.2cm, align=flush left,
          font=\sffamily\footnotesize\color{ink}
        ] {
          \textbf{\small\color{uamred}El problema}\\[3pt]
          674 estudiantes de nuevo ingreso deben quedar
          asignados a grupos del Tronco General (TG) antes
          del inicio del trimestre, respetando: seriación,
          disponibilidad de cupo, franjas horarias sin
          traslape, y prioridad matutina para alumnas.
        };
      \end{tikzpicture}
    \end{column}
    \hfill
    \begin{column}{0.48\textwidth}
      \begin{tikzpicture}
        \node[
          fill=cardwhite, draw=cardborder, line width=0.6pt,
          rounded corners=2pt, inner sep=9pt,
          text width=6.2cm, align=flush left,
          font=\sffamily\footnotesize\color{ink}
        ] {
          \textbf{\small\color{uamred}La solución}\\[3pt]
          Sistema multi-agente con solver de optimización
          combinatoria.\\[3pt]
          Resultado:\\[1pt]
          \textbf{0 conflictos de horario}\\
          \textbf{0 cupos excedidos}\\
          \textbf{3 instancias del solver}
        };
      \end{tikzpicture}
    \end{column}
  \end{columns}

  \vspace{0.25cm}
  \begin{center}
    \kpi{674}{inscritos}%
    \hspace{5pt}%
    \kpi{10}{licenciaturas}%
    \hspace{5pt}%
    \kpi{102}{grupos del TG}%
    \hspace{5pt}%
    \kpi{186 alumnas}{prioridad matutina}%
  \end{center}
\end{frame}
}

%──────────────────────────────────────────────────────────────
% SLIDE 5 — DIVISOR SECCIÓN 02 (fondo oscuro)
%──────────────────────────────────────────────────────────────
{
\setbeamercolor{background canvas}{bg=ink}
\begin{frame}[plain]
  \begin{tikzpicture}[remember picture, overlay]
    \fill[ink] (current page.south west) rectangle (current page.north east);

    \node[anchor=west, text=uamred, opacity=0.18,
          font=\sffamily\bfseries\fontsize{160}{160}\selectfont,
          inner sep=0pt]
      at ($(current page.west)+(-1.0cm,0.0cm)$) {02};

    \draw[uamred, line width=2pt]
      ($(current page.center)+(-3.0cm,1.2cm)$) --
      ($(current page.center)+(-3.0cm,-1.2cm)$);

    \node[anchor=west, text=cardwhite,
          font=\sffamily\bfseries\LARGE, inner sep=0pt]
      at ($(current page.center)+(-2.6cm,0.38cm)$) {Eficiencia terminal};

    \node[anchor=west, text=cardborder,
          font=\sffamily\normalsize, inner sep=0pt]
      at ($(current page.center)+(-2.6cm,-0.38cm)$)
      {Diagnóstico estructural del rezago en los diez planes};
  \end{tikzpicture}
\end{frame}
}

%──────────────────────────────────────────────────────────────
% SLIDE 6 — EFICIENCIA TERMINAL: HALLAZGOS (fondo claro)
%──────────────────────────────────────────────────────────────
{
\setbeamercolor{background canvas}{bg=bg}
\begin{frame}{Eficiencia terminal: hallazgos estructurales}
  \begin{tikzpicture}
    \node[
      fill=cardwhite, draw=cardborder, line width=0.6pt,
      rounded corners=2pt, inner sep=9pt,
      text width=13.0cm, align=flush left,
      font=\sffamily\footnotesize\color{ink}
    ] {
      \textbf{\small\color{uamred}Pareto curricular.}\enspace
      El 26.5\,\% de las seriaciones explica el 80\,\% del beneficio
      potencial de egreso.\\[4pt]
      \textbf{\small\color{uamred}Índices P y h.}\enspace
      Miden la progresión real del estudiante: la brecha promedio entre
      créditos inscritos y la referencia de 38\,cr/trim proviene de las
      cadenas de seriación, no del rendimiento académico.\\[4pt]
      \textbf{\small\color{uamred}Nodos topológicamente críticos.}\enspace
      Un estudiante que reprueba un nodo de este tipo pierde acceso a
      todas las Unidades de Enseñanza-Aprendizaje (UEAs) que dependen de él.
    };
  \end{tikzpicture}

  \vspace{0.3cm}
  \begin{center}
    \kpi{10}{planes analizados}%
    \hspace{5pt}%
    \kpi{5}{infografías interactivas}%
    \hspace{5pt}%
    \kpi{P y h}{índices de progresión}%
  \end{center}
\end{frame}
}

%──────────────────────────────────────────────────────────────
% SLIDE 7 — DIVISOR SECCIÓN 03 (fondo oscuro)
%──────────────────────────────────────────────────────────────
{
\setbeamercolor{background canvas}{bg=ink}
\begin{frame}[plain]
  \begin{tikzpicture}[remember picture, overlay]
    \fill[ink] (current page.south west) rectangle (current page.north east);

    \node[anchor=west, text=uamred, opacity=0.18,
          font=\sffamily\bfseries\fontsize{160}{160}\selectfont,
          inner sep=0pt]
      at ($(current page.west)+(-1.0cm,0.0cm)$) {03};

    \draw[uamred, line width=2pt]
      ($(current page.center)+(-3.0cm,1.2cm)$) --
      ($(current page.center)+(-3.0cm,-1.2cm)$);

    \node[anchor=west, text=cardwhite,
          font=\sffamily\bfseries\LARGE, inner sep=0pt]
      at ($(current page.center)+(-2.6cm,0.38cm)$) {Distribución docente};

    \node[anchor=west, text=cardborder,
          font=\sffamily\normalsize, inner sep=0pt]
      at ($(current page.center)+(-2.6cm,-0.38cm)$)
      {Concentración y riesgo curricular 2016--2025};
  \end{tikzpicture}
\end{frame}
}

%──────────────────────────────────────────────────────────────
% SLIDE 8 — DISTRIBUCIÓN DOCENTE: HALLAZGOS (fondo claro)
%──────────────────────────────────────────────────────────────
{
\setbeamercolor{background canvas}{bg=bg}
\begin{frame}[t]{Distribución docente: cuatro hallazgos}
  \begin{tikzpicture}
    \node[
      fill=cardwhite, draw=cardborder, line width=0.6pt,
      rounded corners=2pt, inner sep=6pt,
      text width=13.0cm, align=flush left,
      font=\sffamily\footnotesize\color{ink}
    ] {
      \textbf{\small\color{uamred}Monopolios de impartición.}\enspace
      187 UEAs (23.1\,\%) tienen HHI = 1:
      un único profesor ha impartido esa UEA durante toda su historia registrada.\\[2pt]
      \textbf{\small\color{uamred}Riesgo crítico.}\enspace
      61 UEAs en riesgo crítico (RC $\geq$ percentil~90): alta concentración
      en posición estructuralmente relevante de la cadena de seriación.\\[2pt]
      \textbf{\small\color{uamred}Fragilidad sistémica.}\enspace
      11 de 22 vecindades curriculares detectadas presentan fragilidad
      sistémica, no puntual.\\[2pt]
      \textbf{\small\color{uamred}Paradoja de la concentración.}\enspace
      El cluster de mayor HHI agregado tiene cero UEAs críticas;
      el de menor HHI tiene 19:
      la posición en la cadena importa más que la concentración.
    };
  \end{tikzpicture}

  \vspace{0.1cm}
  \begin{center}
    \kpi{809}{UEAs activas}%
    \hspace{5pt}%
    \kpi{791}{profesores}%
    \hspace{5pt}%
    \kpi{5{,}391 grupos}{2016--2025}%
  \end{center}
\end{frame}
}

%──────────────────────────────────────────────────────────────
% SLIDE 9 — DIVISOR SECCIÓN 04 (fondo oscuro)
%──────────────────────────────────────────────────────────────
{
\setbeamercolor{background canvas}{bg=ink}
\begin{frame}[plain]
  \begin{tikzpicture}[remember picture, overlay]
    \fill[ink] (current page.south west) rectangle (current page.north east);

    \node[anchor=west, text=uamred, opacity=0.18,
          font=\sffamily\bfseries\fontsize{160}{160}\selectfont,
          inner sep=0pt]
      at ($(current page.west)+(-1.0cm,0.0cm)$) {04};

    \draw[uamred, line width=2pt]
      ($(current page.center)+(-3.0cm,1.2cm)$) --
      ($(current page.center)+(-3.0cm,-1.2cm)$);

    \node[anchor=west, text=cardwhite,
          font=\sffamily\bfseries\LARGE, inner sep=0pt]
      at ($(current page.center)+(-2.6cm,0.38cm)$) {Modificaciones 2026};

    \node[anchor=west, text=cardborder,
          font=\sffamily\normalsize, inner sep=0pt]
      at ($(current page.center)+(-2.6cm,-0.38cm)$)
      {Similitud semántica y revisión legislativa automatizada};
  \end{tikzpicture}
\end{frame}
}

%──────────────────────────────────────────────────────────────
% SLIDE 10 — MODIFICACIONES 2026: DOS SISTEMAS (fondo claro)
%──────────────────────────────────────────────────────────────
{
\setbeamercolor{background canvas}{bg=bg}
\begin{frame}[t]{Modificaciones 2026: dos sistemas de apoyo}
  \begin{columns}[T]
    \begin{column}{0.48\textwidth}
      \begin{tikzpicture}
        \node[
          fill=cardwhite, draw=cardborder, line width=0.6pt,
          rounded corners=2pt, inner sep=7pt,
          text width=6.2cm, align=flush left,
          font=\sffamily\footnotesize\color{ink}
        ] {
          \textbf{\small\color{uamred}Análisis de similitud}\\[2pt]
          1{,}505 pares de UEAs comparados con inteligencia
          artificial entre los diez planes.\\[1pt]
          203 pares idénticos o casi idénticos.\\[1pt]
          El par Industrial--Mecánica tiene 38 programas
          equivalentes: candidatos a UEAs transversales
          compartidas.
        };
      \end{tikzpicture}
    \end{column}
    \hfill
    \begin{column}{0.48\textwidth}
      \begin{tikzpicture}
        \node[
          fill=cardwhite, draw=cardborder, line width=0.6pt,
          rounded corners=2pt, inner sep=7pt,
          text width=6.2cm, align=flush left,
          font=\sffamily\footnotesize\color{ink}
        ] {
          \textbf{\small\color{uamred}Revisión legislativa}\\[2pt]
          5 agentes encadenados verifican cada propuesta contra
          el Reglamento de Estudios Superiores (RES):
          10~fracciones del Art.~34, 16 del Art.~44, más
          Art.~45 y~56.\\[1pt]
          Salida: 2~reportes Excel con veredicto
          (Sí\,/\,Parcial\,/\,No\,/\,N\!A)
          por fracción.\\[1pt]
          10 oficios generados.
        };
      \end{tikzpicture}
    \end{column}
  \end{columns}

  \vspace{0.1cm}
  \begin{center}
    \kpi{1{,}505 pares}{de UEAs}%
    \hspace{5pt}%
    \kpi{626}{programas}%
    \hspace{5pt}%
    \kpi{26 fracciones}{del RES}%
    \hspace{5pt}%
    \kpi{10 oficios}{generados}%
  \end{center}
\end{frame}
}

%──────────────────────────────────────────────────────────────
% SLIDE 11 — CIERRE / ACCESO (fondo oscuro)
%──────────────────────────────────────────────────────────────
{
\setbeamercolor{background canvas}{bg=ink}
\begin{frame}[plain]
  \begin{tikzpicture}[remember picture, overlay]
    \fill[ink] (current page.south west) rectangle (current page.north east);

    % Texto principal
    \node[anchor=north, text=cardwhite,
          font=\sffamily\bfseries\Large,
          inner sep=0pt]
      at ($(current page.north)+(0cm,-1.2cm)$) {
        El dashboard está disponible en línea
    };

    % QR grande centrado
    \node[anchor=north, inner sep=0pt]
      at ($(current page.north)+(0cm,-2.2cm)$) {
        \includegraphics[width=4.0cm]{%
          /Users/nowhere/Claude_Projects/Admin_duties/dashboard/figuras/qr_dashboard_dcbi.png}
    };

    % URL — debajo del QR
    \node[anchor=north, text=uamred,
          font=\sffamily\normalsize\bfseries, inner sep=0pt]
      at ($(current page.north)+(0cm,-6.5cm)$) {
        cslmuam.github.io/dashboard-dcbi-2026
    };

    % Línea separadora
    \draw[uamred, line width=0.6pt]
      ($(current page.south west)+(0.6cm,2.0cm)$) --
      ($(current page.south east)+(-0.6cm,2.0cm)$);

    % Contacto — debajo de la línea
    \node[anchor=south, text=cardborder,
          font=\sffamily\footnotesize, inner sep=0pt]
      at ($(current page.south)+(0cm,1.1cm)$) {
        Dr.\ César Simón López-Monsalvo
        \enspace\textbf{\color{uamred}·}\enspace
        cslm@azc.uam.mx
%        \enspace\textbf{\color{uamred}·}\enspace
%        Departamento de Ciencias Básicas
    };

  \end{tikzpicture}
\end{frame}
}

\end{document}
